\chapter{Diseño experimental}\label{cap:diseno-experimental}

% Depende de H2 (protocolo pre-registrado). Fuente: "Marco metodológico" de la propuesta.

\section{Marco metodológico}
% Design Science Research (Hevner et al. 2004) + estudio de caso (Yin 2018).

\section{Diseño factorial 2×2}
% Factores: modelo (A/B) × RAG (sin/con); las cuatro celdas.

\section{Variables dependientes}
% Tasa de ATs superados (holdout), intervenciones por causa raíz, alucinaciones de
% dominio, métricas estáticas, adherencia a estándares; rol desarrollador vs. revisor.

\section{Variables controladas y protocolo de corrida}
% Spec pinneada, corpus congelado, evaluador único, criterios de intervención
% pre-registrados, presupuestos, orden de construcción.

\section{Amenazas a la validez}
% n=1 por celda (sin inferencia estadística), deriva de modelos comerciales, sesgo
% del evaluador (mitigado por harness pre-construido), generalización del caso único.
